\chapter{What Dopamine Really Says}
% full version: chapters/07_dofaalt.tex

In the last chapter dopamine was simple: one number, ``better or worse than expected.'' That is a first approximation, and it works well. But experiments in recent years have shown that more than one number runs along this ``wire.'' This chapter is about where the picture gets more complicated and where scientists are still arguing.

\section*{Optimists and Pessimists}

Dopamine neurons are not all alike. Some respond more strongly to pleasant surprises than to unpleasant ones: these are optimists. Others do the opposite: pessimists. If each neuron learns by its own rule, the optimist learns what the reward will be at best, and the pessimist at worst. Together they record not one average expectation but the whole spread of possible outcomes, like a bookmaker who knows not only the favorite but everyone's odds. The idea agrees with experiment; why the brain needs the whole spread is still a hypothesis.

\pencilfig{7}{\pencilart{7}}{Each dopamine neuron holds its own point on the curve of possible rewards: together they remember the whole curve.}

\section*{Not Only ``Better'' but Also ``Important''}

Some dopamine neurons also flare up at unpleasant but important events: a loud sound, danger. It seems the dopamine network carries not only the sign of the surprise but also its significance, a ``pay attention'' signal.

\section*{Two Signals on One Wire}

The fast bursts teach. But the slow average level of dopamine, which changes over minutes, seems to say something else: how rich the surroundings are right now. If there is plenty of reward around, time is precious, and the animal acts faster. Experiments confirm that the dopamine level is linked to the vigor of actions. An engineer would call this separating signals by frequency: the high frequencies carry one thing, the low frequencies another.

\section*{A Ramp Near the Goal}

When an animal approaches a reward, dopamine often does not flare but rises smoothly. One explanation: this is the expectation itself, a signal saying ``the goal is close, push on.'' Another keeps the surprise idea: from far away the animal judges its surroundings vaguely, more and more precisely as it gets closer, and each refinement is a small pleasant surprise. An elegant experiment in a virtual corridor, where the animal was instantly moved closer to the goal, favors the second: dopamine responded with a burst, as a surprise should.

\section*{Looking Back}

There is also a bold alternative: dopamine answers not the question ``what will happen?'' but the question ``what caused the reward?'' It looks into the past rather than the future. Experiments in which the two theories predict different things contradict one another so far. This is a live scientific debate.

\section*{A Vector Instead of a Number}

Finally, dopamine also responds to a change in the \emph{flavor} of the reward when its value has not changed, and it helps learning without any reward at all. It seems this is not one error but a whole bundle of errors, one for each feature of what happened.

The chapter's conclusion: most of the new theories do not overturn the simple picture but extend it. Dopamine is not one wire but a small bundle of several. Only ``looking back'' seriously disputes it.

\fullversion{7}
